\documentclass[12pt,a4paper]{article}
\usepackage[utf8]{inputenc}
\usepackage[T1]{fontenc}
\usepackage[french]{babel}
\usepackage{amsmath, amssymb, amsthm, mathrsfs}
\usepackage{geometry}
\usepackage{newunicodechar}
\newunicodechar{∀}{\ensuremath{\forall}}
\newunicodechar{∃}{\ensuremath{\exists}}
\newunicodechar{≥}{\ensuremath{\ge}}
\newunicodechar{≤}{\ensuremath{\le}}
\newunicodechar{∈}{\ensuremath{\in}}
\newunicodechar{×}{\ensuremath{\times}}
\newunicodechar{∧}{\ensuremath{\land}}
\newunicodechar{∩}{\ensuremath{\cap}}
\newunicodechar{ℂ}{\ensuremath{\mathbb{C}}}
\newunicodechar{ℝ}{\ensuremath{\mathbb{R}}}
\newunicodechar{∑}{\ensuremath{\sum}}
\newunicodechar{∫}{\ensuremath{\int}}
\newunicodechar{‖}{\ensuremath{\|}}
\newunicodechar{θ}{\ensuremath{\theta}}

\geometry{margin=2.5cm}

\newtheorem{theorem}{Th\'eor\`eme}[section]
\newtheorem{lemma}[theorem]{Lemme}
\newtheorem{definition}[theorem]{D\'efinition}
\newtheorem{conjecture}[theorem]{Conjecture}
\newtheorem{corollary}[theorem]{Corollaire}
\newtheorem{proposition}[theorem]{Proposition}

\title{Sur la Conjecture d'Erd\H{o}s-Tur\'an pour les Bases Additives :\\ Analyse Harmonique Discr\`ete, M\'ethode du Cercle et Formalisation Topologique}
\author{D\'epartement de Math\'ematiques Pures}
\date{}

\begin{document}
\maketitle
\tableofcontents

\section{Introduction G\'en\'erale et Cadre Axiomatique}

La th\'eorie additive des nombres s'attache \`a la compr\'ehension des propri\'et\'es de repr\'esentation des entiers par des sous-ensembles s\'electionn\'es. Parmi les d\'efis historiques de cette discipline, la conjecture d'Erd\H{o}s-Tur\'an, non r\'esolue depuis 1941, \'etablit une fronti\`ere in\'evitable sur la r\'egularit\'e asymptotique des d\'ecompositions arithm\'etiques. L'objet de cette monographie est d'\'etablir les structures axiomatiques formelles n\'ecessaires \`a la compr\'ehension fine de ce ph\'enom\`ene de fluctuation in\'evitable.

\subsection{D\'efinitions Pr\'eliminaires et Typage Rigoureux}

Soit l'ensemble des entiers naturels $\mathbb{N} = \{0, 1, 2, \dots\}$. Toutes les variables d\'enotant des ensembles (usuellement $A, B$) appartiennent \`a $\mathcal{P}(\mathbb{N})$, l'ensemble des parties de $\mathbb{N}$. Les variables indices (usuellement $n, m, k, a, b$) appartiennent \`a $\mathbb{N}$.

Nous d\'efinissons la notion de base additive d'ordre 2 de mani\`ere stricte.

\begin{definition}
Un sous-ensemble $B \subseteq \mathbb{N}$ est d\'enomm\'e \textit{base additive d'ordre $2$ asymptotique} si le compl\'ementaire de l'ensemble des entiers repr\'esentables comme la somme de deux \'el\'ements de $B$ est fini. Formellement, $B \in \mathcal{P}(\mathbb{N})$ satisfait la condition :
\begin{equation}
\exists N_0 \in \mathbb{N}, \quad \forall n \in \mathbb{N}, \quad n \ge N_0 \implies \exists a \in B, \exists b \in B, a + b = n.
\end{equation}
\end{definition}

La quantit\'e des diff\'erentes paires ordonn\'ees $(a, b)$ solutions pour un m\^eme entier $n$ d\'etermine la structure de superposition des sommes.

\begin{definition}
La fonction de repr\'esentation associ\'ee \`a l'ensemble $B \subseteq \mathbb{N}$ et \`a l'ordre 2 est l'application $r_{B,2} : \mathbb{N} \to \mathbb{N}$ d\'efinie par le cardinal de l'ensemble des couples valides :
\begin{equation}
r_{B,2}(n) = \left| \{ (a, b) \in B \times B \mid a + b = n \} \right|.
\end{equation}
\end{definition}

Par construction, l'ensemble $S_n = \{ (a, b) \in B \times B \mid a + b = n \}$ est inclus dans l'ensemble fini $\{ (a, b) \in \mathbb{N} \times \mathbb{N} \mid a + b = n \}$, dont le cardinal est exactement $n + 1$. Par cons\'equent, l'ensemble $S_n$ est n\'ecessairement fini, garantissant que $r_{B,2}(n)$ est une fonction bien d\'efinie \`a valeurs dans $\mathbb{N}$.

La condition pour que $B$ soit une base asymptotique d'ordre 2 se r\'e\'ecrit donc sous la forme unifi\'ee $r_{B,2}(n) \ge 1$ pour tout entier $n \ge N_0$.

\begin{conjecture}[Erd\H{o}s-Tur\'an, 1941]
Soit $B \subseteq \mathbb{N}$. Si $r_{B,2}(n) > 0$ pour tout entier $n$ au-del\`a d'un certain seuil $N_0$, alors la limite sup\'erieure de la s\'uite $(r_{B,2}(n))_{n \in \mathbb{N}}$ est infinie :
\begin{equation}
\limsup_{n \to \infty} r_{B,2}(n) = \infty.
\end{equation}
\end{conjecture}

L'argumentaire de l'approche analytique d\'evelopp\'ee dans ce texte proc\`ede par d\'emonstration par l'absurde. Nous formulons donc l'hypoth\`ese n\'egative stricte.

\begin{definition}
L'hypoth\`ese de r\'egularit\'e globale born\'ee (HRA) consiste \`a postuler l'existence d'une constante globale ind\'ependante $M \in \mathbb{N}$ telle que :
\begin{equation}
\forall n \in \mathbb{N}, \quad r_{B,2}(n) \le M.
\end{equation}
\end{definition}

Dans la d\'efinition de HRA, il est clair que si $r_{B,2}(n) \le M$ est vrai pour tout $n \ge N_0$, comme $r_{B,2}(n)$ est fini pour tout $n < N_0$, on peut r\'eajuster $M$ en prenant le maximum de $M$ et de $\max_{0 \le k < N_0} r_{B,2}(k)$ pour obtenir une borne globale valable pour tout $n \ge 0$. C'est cette forme forte que nous utiliserons sans perte de g\'en\'eralit\'e.

\section{Litt\'erature Contextuelle et Analogies Structurelles}

L'\'etude de la fonction $r_{B,2}(n)$ se situe au confluent de la combinatoire, de la th\'eorie analytique des nombres et de l'analyse harmonique discr\`ete. Historiquement, l'\'etude des bases additives a d\'ebut\'e avec les th\'eor\`emes de Lagrange sur les quatre carr\'es et les probl\`emes de Waring, o\`u les bases pr\'esentent une structure alg\'ebrique rigide. Cependant, la conjecture d'Erd\H{o}s-Tur\'an interroge la nature intrins\`eque de \textit{toute} base additive, ind\'ependamment de sa d\'efinition alg\'ebrique.

\subsection{Th\'eor\`emes G\'en\'eraux de Fluctuation}

Le premier jalon significatif dans la direction de la compr\'ehension des fluctuations n\'ecessaires de $r_{B,2}(n)$ fut pos\'e par P.A.M. Dirac en 1951. Dirac a d\'emontr\'e que si une s\'uite $B$ poss\`ede une densit\'e asymtotique trop r\'eguli\`ere, la fonction de repr\'esentation doit osciller. Toutefois, la caract\'erisation la plus puissante des limites de la r\'egularit\'e asymptotique fut apport\'ee par le th\'eor\`eme d'Erd\H{o}s-Fuchs en 1956.

Le th\'eor\`eme d'Erd\H{o}s-Fuchs d\'emontre qu'il est analytiquement impossible pour la somme partielle des fonctions de repr\'esentation de s'aligner sur une progression lin\'eaire parfaite avec un terme d'erreur n\'egligeable. Formellement, il n'existe aucune constante $c > 0$ telle que :
\begin{equation}
\sum_{k=0}^n r_{B,2}(k) = c \cdot n + o\left(n^{1/4} (\log n)^{-1/2}\right).
\end{equation}
La d\'emonstration de ce th\'eor\`eme repose sur la m\'ethode du cercle et l'int\'egration complexe de la fonction g\'en\'eratrice. Si $r_{B,2}(n)$ \'etait born\'e (comme le postule notre HRA), ses sommes partielles seraient lin\'eairement born\'ees, mais la variance de cette somme partielle autour de sa moyenne ne pourrait \^etre r\'eduite au-del\`a de la borne d'Erd\H{o}s-Fuchs. Bien que le th\'eor\`eme d'Erd\H{o}s-Fuchs n'interdise pas formellement que $r_{B,2}(n)$ reste constant (et donc borne), il prouve qu'aucune "lisse" distribution n'est possible, sugg\'erant fortement que l'invariance est impossible.

\subsection{Analogie avec le Th\'eor\`eme de Szemer\'edi et Green-Tao}

Une analogie structurelle et m\'ethodologique profonde existe entre le probl\`eme d'Erd\H{o}s-Tur\'an et les d\'eveloppements r\'ecents sur l'existence de progressions arithm\'etiques dans les sous-ensembles denses des entiers. Le Th\'eor\`eme de Roth (1953) \'etablit que tout sous-ensemble des entiers de densit\'e sup\'erieure positive contient des progressions arithm\'etiques de longueur 3. La preuve de Roth utilise l'analyse de Fourier discr\`ete pour d\'etecter les biais dans la distribution modulo un certain entier.

La s\'eparation formelle entre un comportement "structur\'e" (forts coefficients de Fourier) et "pseudo-al\'eatoire" (faibles coefficients de Fourier uniform\'ement r\'epartis), conceptualis\'ee plus tard par Gowers (normes d'uniformit\'e) et aboutissant au th\'eor\`eme de Green-Tao (les nombres premiers contient des progressions arithm\'etiques de longueur arbitraire), repose sur les m\^emes objets analytiques : les sommes d'exponentielles.

Dans le cas d'Erd\H{o}s-Tur\'an, l'ensemble $B$ doit \^etre "dense" (au sens de la fonction de comptage $B(N) \ge c \sqrt{N}$ que nous d\'emontrerons) pour \^etre une base. Cette densit\'e garantit, \`a l'instar du th\'eor\`eme de Roth, un "gros" coefficient de Fourier \`a l'origine. L'hypoth\`ese de r\'egularit\'e globale $r_{B,2}(n) \le M$ force, par l'identit\'e de Parseval, la norme $L^4$ du polyn\^ome de Fourier \`a \^etre fortement contrainte. La contradiction analytique \'emmerge du fait qu'une fonction Lipschitzienne poss\'edant une masse sp\'ecifique \`a l'origine ne peut pas \'ecraser sa norme $L^4$ int\'egrale de mani\`ere uniforme sans violer l'in\'egalit\'e isop\'erim\'etrique discr\`ete inf\'er\'ee par le th\'eor\`eme des accroissements finis sur le cercle unitaire. Les outils spectraux sont identiques \`a ceux employ\'es dans l'\'etude des progressions arithm\'etiques.

\section{M\'ecanismes G\'en\'erateurs et \'Equivalence Fonctionnelle Alg\'ebrique}

L'\'etude de $r_{B,2}(n)$ d\'epend intrins\`equement des fonctions g\'en\'eratrices complexes. Pour \'eviter les probl\`emes de convergence globale, nous utilisons des polyn\^omes de degr\'e fini (s\'eries de Taylor tronqu\'ees).

\subsection{Construction de la S\'erie Restreinte}

Soit un nombre entier strict $N \in \mathbb{N}$, $N \ge 1$. On construit un polyn\^ome alg\'ebrique complexe dont le support est la restriction de la base additive \`a l'intervalle $[0, N]$. Pour une variable $z \in \mathbb{C}$ d\'efinie sur le disque unit\'e ferm\'e $\{ z \in \mathbb{C} \mid |z| \le 1 \}$, la fonction indicatrice g\'en\'eratrice se d\'efinit formellement par :
\begin{equation}
P_N(z) = \sum_{\substack{b \in B \\ 0 \le b \le N}} z^b.
\end{equation}

Puisque l'ensemble de sommation est fini, l'\'evaluation analytique de $P_N(z)$ est bien d\'efinie et la fonction est ind\'efiniment diff\'erentiable (enti\`ere) sur $\mathbb{C}$.

Le th\'eor\`eme de convolution fondamental pour les sommes directes finies assure la propri\'et\'e morphique g\'en\'eratrice. En d\'eveloppant analytiquement $P_N(z)^2$, l'expansion des puissances s'agr\`ege selon les invariants combinatoires.

\begin{lemma}[D\'eveloppement Quadratique Restreint]
Soit $r_{\le N}(n)$ l'application d\'efinie sur $\mathbb{N}$ par le cardinal de l'ensemble des paires $(a,b) \in B \times B$ telles que $a \le N, b \le N$ et $a+b = n$.
Pour tout $z \in \mathbb{C}$ :
\begin{equation}
P_N(z)^2 = \sum_{n=0}^{2N} r_{\le N}(n) z^n.
\end{equation}
De plus, par inclusion stricte des ensembles supports, on a inconditionnellement pour tout entier $n \in \mathbb{N}$ l'in\'egalit\'e :
\begin{equation}
r_{\le N}(n) \le r_{B,2}(n).
\end{equation}
\end{lemma}

\begin{proof}
L'identit\'e est formellement alg\'ebrique. Par d\'efinition de la multiplication des sommes finies :
\begin{equation}
P_N(z)^2 = \left( \sum_{\substack{a \in B \\ 0 \le a \le N}} z^a \right) \left( \sum_{\substack{b \in B \\ 0 \le b \le N}} z^b \right).
\end{equation}
En appliquant la distributivit\'e de la multiplication sur l'addition dans le corps $\mathbb{C}$, on obtient une somme double :
\begin{equation}
P_N(z)^2 = \sum_{\substack{a \in B \\ 0 \le a \le N}} \sum_{\substack{b \in B \\ 0 \le b \le N}} z^a z^b = \sum_{\substack{a \in B \\ 0 \le a \le N}} \sum_{\substack{b \in B \\ 0 \le b \le N}} z^{a+b}.
\end{equation}
Puisque $0 \le a \le N$ et $0 \le b \le N$, la somme $a+b$ est n\'ecessairement encadr\'ee par $0 \le a+b \le 2N$. Nous partitionnons l'ensemble des paires valides d'indices $(a,b)$ selon la valeur de leur somme, not\'ee $n$. L'ensemble de toutes les paires possibles $(a,b) \in B \cap [0, N] \times B \cap [0, N]$ se r\'epartit dans l'union disjointe des ensembles $E_n = \{ (a,b) \in B \times B \mid a \le N \land b \le N \land a+b = n \}$ pour $n \in [0, 2N]$.
Le cardinal de $E_n$ est pr\'ecis\'ement la d\'efinition donn\'ee pour $r_{\le N}(n)$. La somme double se r\'ecrit alors exactement en regroupant les termes ayant le m\^eme exposant $n$ :
\begin{equation}
\sum_{\substack{a \in B \\ 0 \le a \le N}} \sum_{\substack{b \in B \\ 0 \le b \le N}} z^{a+b} = \sum_{n=0}^{2N} \left( \sum_{(a,b) \in E_n} z^n \right) = \sum_{n=0}^{2N} |E_n| z^n = \sum_{n=0}^{2N} r_{\le N}(n) z^n.
\end{equation}
Ceci conclut la d\'emonstration de l'\'egalit\'e alg\'ebrique analytique.

Pour d\'emontrer l'in\'egalit\'e $r_{\le N}(n) \le r_{B,2}(n)$, il suffit d'utiliser la monotonie de la mesure de comptage sur les ensembles finis. Par d\'efinition :
\begin{equation}
E_n = \{ (a,b) \in B \times B \mid a \le N \land b \le N \land a+b = n \}.
\end{equation}
Et :
\begin{equation}
S_n = \{ (a,b) \in B \times B \mid a+b = n \}.
\end{equation}
La condition $a \le N \land b \le N$ est restrictive. Si une paire $(a,b)$ appartient \`a $E_n$, elle appartient \textit{a fortiori} \`a $S_n$. Donc $E_n \subseteq S_n$. Puisque $S_n$ est un ensemble fini, l'inclusion des sous-ensembles implique l'in\'egalit\'e sur les cardinaux, soit $|E_n| \le |S_n|$, ce qui s'\'ecrit strictement $r_{\le N}(n) \le r_{B,2}(n)$. La d\'emonstration du lemme est rigoureusement \'etablie.
\end{proof}

\subsection{D\'erivation de la Borne Inf\'erieure de Croissance Lin\'eaire}

L'exploitation de la propri\'et\'e d\'efinissant $B$ comme base asymptotique d'ordre 2, soit $r_{B,2}(n) \ge 1$ pour tout $n \ge N_0$, permet d'imposer formellement une masse minimale aux sommes g\'en\'eratrices.

\begin{lemma}[Estimation Lin\'eaire de Masse D\'eduite]
Sous l'hypoth\`ese de base additive asymptotique, il existe une constante positive r\'eelle $C$, ind\'ependante de $N$, telle que pour tout $N > N_0$ :
\begin{equation}
\sum_{n=0}^N r_{\le N}(n) \ge N - C.
\end{equation}
\end{lemma}

\begin{proof}
L'\'el\'ement central de la d\'emonstration consiste \`a observer l'intersection exacte des conditions de restriction.
Fixons un entier arbitraire $n$ tel que $0 \le n \le N$.
Consid\'erons l'ensemble $S_n = \{ (a,b) \in B \times B \mid a+b = n \}$. Soit $(a,b)$ un \'el\'ement quelconque appartenant \`a $S_n$. Puisque les entiers appartenant \`a l'ensemble $B$ sont des entiers naturels, ils sont n\'ecessairement positifs ou nuls. Ainsi, $a \ge 0$ et $b \ge 0$.
L'\'equation $a+b=n$ implique, par soustraction d'entiers positifs, que $a = n - b \le n$. Puisque par choix initial de l'indexation $n \le N$, la transitivit\'e de l'ordre usuel sur $\mathbb{N}$ dicte que $a \le N$. Par le m\^eme argument analytique de sym\'etrie, $b = n - a \le n \le N$.
Par cons\'equent, toute paire $(a,b)$ r\'esolvant $a+b=n$ avec $n \le N$ satisfait syst\'ematiquement les deux conditions restrictives $a \le N$ et $b \le N$.
Cela entra\^ine formellement l'\'egalit\'e des ensembles $S_n$ et $E_n$ pour tout $n \in [0, N]$. L'\'egalit\'e de cardinaux en d\'ecoule imm\'ediatement :
\begin{equation}
r_{\le N}(n) = r_{B,2}(n), \quad \forall n \in [0, N].
\end{equation}

Consid\'erons la somme partielle d\'efinie dans l'\'enonc\'e du Lemme et ins\'erons cette substitution exacte :
\begin{equation}
\sum_{n=0}^N r_{\le N}(n) = \sum_{n=0}^N r_{B,2}(n).
\end{equation}
Divisons l'intervalle d'int\'egration discr\`ete en la singularit\'e temporelle de la d\'efinition, $N_0$. Nous supposons validement $N > N_0$ par hypoth\`ese.
\begin{equation}
\sum_{n=0}^N r_{B,2}(n) = \sum_{n=0}^{N_0 - 1} r_{B,2}(n) + \sum_{n=N_0}^N r_{B,2}(n).
\end{equation}
Posons la valeur globale constante repr\'esentant les fluctuations r\'esiduelles locales : $\gamma = \sum_{n=0}^{N_0 - 1} r_{B,2}(n)$. L'ensemble $B$ \'etant fix\'e et ind\'ependant de $N$, le param\`etre $N_0$ est sp\'ecifique \`a l'ensemble abstrait $B$ ind\'ependamment de la coupure $N$. Ainsi $\gamma$ est une constante finie absolue.
Pour le second terme de la partition, l'application it\'er\'ee de l'axiome $r_{B,2}(n) \ge 1$ sur le support de sommation produit la minoration analytique stricte :
\begin{equation}
\sum_{n=N_0}^N r_{B,2}(n) \ge \sum_{n=N_0}^N 1.
\end{equation}
L'\'evaluation de cette somme de termes unitaires r\'esulte en le comptage de la longueur de l'intervalle discret : $(N - N_0 + 1)$.
En r\'e-assemblant les deux intervalles, nous obtenons l'in\'egalit\'e lin\'eaire globale :
\begin{equation}
\sum_{n=0}^N r_{\le N}(n) \ge \gamma + (N - N_0 + 1).
\end{equation}
En restructurant alg\'ebriquement le membre de droite pour isoler la composante proportionnelle $N$, et en d\'efinissant formellement la constante $C = N_0 - 1 - \gamma$, nous concluons de fa\c{c}on rigoureuse :
\begin{equation}
\sum_{n=0}^N r_{\le N}(n) \ge N - C.
\end{equation}
La preuve est aboutie, les d\'eductions encha\^in\'ees d\'emontrent logiquement la minoration asymptotique en $\mathcal{O}(N)$ sur la restriction partielle de l'op\'erateur de densit\'e.
\end{proof}

Cette in\'egalit\'e analytique autorise directement la formation d'un theoreme sur la croissance asymtotique de l'ensemble racine $B$ lui-m\^eme.

\begin{theorem}[Borne Inf\'erieure Stricte de la Densit\'e de la Base]
Soit l'op\'erateur de d\'enombrement continu discret $B(N)$ d\'efini par le cardinal $|B \cap [0, N]|$. Sous l'hypoth\`ese que $B$ est une base additive d'ordre 2, le terme source cro\^it asymptotiquement d'au moins la m\'etrique racine :
\begin{equation}
B(N) \ge \sqrt{N - C},
\end{equation}
pour tout $N > N_0$.
\end{theorem}

\begin{proof}
L'objet focal est la formulation restreinte de la g\'en\'eratrice complexe sur le point fondamental $z=1 \in \mathbb{C}$.
L'\'evaluation scalaire se lit par d\'efinition :
\begin{equation}
P_N(1) = \sum_{\substack{b \in B \\ 0 \le b \le N}} 1^b = \sum_{\substack{b \in B \\ 0 \le b \le N}} 1 = | B \cap [0, N] | = B(N).
\end{equation}
Dans un deuxi\`eme temps, l'\'evaluation analytique du carr\'e de l'indicatrice $P_N(z)$ en $z=1$ au travers du d\'eveloppement polynomial formel du Lemme 3.1 donne :
\begin{equation}
P_N(1)^2 = \left( \sum_{n=0}^{2N} r_{\le N}(n) 1^n \right) = \sum_{n=0}^{2N} r_{\le N}(n).
\end{equation}
Ainsi, le d\'enombrement scalaire quadratique $B(N)^2$ caract\'erise de mani\`ere univoque l'int\'egralit\'e des repr\'esentations restreintes.
Sachant formellement par construction des espaces cardinaux finis que pour tout sous-indice temporel $n$, on a n\'ecessairement $r_{\le N}(n) \ge 0$, la troncature sup\'erieure d'une s\'erie finie n'induit qu'une r\'eduction strictement positive ou nulle :
\begin{equation}
\sum_{n=0}^{2N} r_{\le N}(n) \ge \sum_{n=0}^{N} r_{\le N}(n).
\end{equation}
En substituant le r\'esultat de la majoration de la borne inf\'erieure \'etablie explicitement dans le Lemme 3.2 sur l'intervalle restreint, l'in\'egalit\'e devient imm\'ediate :
\begin{equation}
P_N(1)^2 \ge N - C.
\end{equation}
Il est donn\'e que l'application analytique racine carr\'ee continue, soit $f(x) = \sqrt{x}$, est un hom\'eomorphisme strictement croissant et pr\'eserve de facto les in\'egalit\'es de r\'eels positifs. Appliqu\'ee sur le syst\`eme in\'egalitaire d\'emontr\'e (le membre de droite \'etant asymptotiquement grand et positif), l'op\'eration restitue :
\begin{equation}
P_N(1) \ge \sqrt{N - C},
\end{equation}
ce qui confirme l'in\'egalit\'e annonc\'ee.
\end{proof}

\section{Formalisation Duale et In\'egalit\'es Spectrales Int\'egrales}

La m\'ethode du cercle de Hardy-Littlewood se caract\'erise par l'utilisation de l'int\'egration le long de sous-espaces r\'eguliers dans le domaine des composantes complexes. Le domaine r\'ealisable est pr\'ecis\'ement le cercle unit\'e topologique d\'efini num\'eriquement par $\theta \mapsto e^{i\theta}$ sur l'intervalle modulaire $\theta \in [0, 2\pi)$.

\subsection{Identit\'e Spectrale Isom\'etrique par Produit Scalaire}

La premi\`ere application est l'obtention d'une \'equivalence \'energ\'etique, g\'en\'eralisant l'identit\'e de Parseval-Plancherel pour les fonctions discr\`etes enti\`eres. Le lemme suivant r\'ealise l'orthonormalisation exacte.

\begin{lemma}[Th\'eor\`eme Orthonormal de l'\'Energie de Parseval]
\label{lem:parseval}
Pour tout param\`etre d'\'epaisseur temporelle $N > 0$, l'\'equation int\'egro-diff\'erentielle des projections conjugu\'ees de phase se formalise exactement par :
\begin{equation}
\frac{1}{2\pi} \int_0^{2\pi} |P_N(e^{i\theta})|^4 d\theta = \sum_{n=0}^{2N} r_{\le N}(n)^2.
\end{equation}
\end{lemma}

\begin{proof}
La propri\'et\'e du module quadratique sur l'espace int\'egre complexe $|z|^4 = (z \overline{z})^2 = z^2 \overline{z^2}$ est exploit\'ee au premier rang analytique. L'\'evaluation de $P_N(z)$ restreinte au contour ferm\'e s'identifie pour un argument quelconque $\theta$. Le r\'esultat morphique d\'emontr\'e de la fonction g\'en\'eratrice nous permet de formuler $P_N(e^{i\theta})^2 = \sum_{n=0}^{2N} r_{\le N}(n) e^{in\theta}$.
L'application de la conjugaison complexe sur une s\'erie de param\`etres absolus r\'eels g\'en\`ere la distribution sym\'etrique de phase :
\begin{equation}
\overline{P_N(e^{i\theta})^2} = \overline{ \sum_{n=0}^{2N} r_{\le N}(n) e^{in\theta} } = \sum_{n=0}^{2N} r_{\le N}(n) \overline{e^{in\theta}} = \sum_{m=0}^{2N} r_{\le N}(m) e^{-im\theta}.
\end{equation}
Ici nous renommons l'indice op\'eratoire de l'espace conjugu\'e de $n$ en $m$ pour l'\'evitement de collision s\'emantique lors des amalgames it\'eratifs, afin d'observer formellement le d\'eveloppement.
Le d\'eploiement du module en s\'erie entrelac\'ee produit une structure double par multiplicativit\'e de la sommation :
\begin{equation}
|P_N(e^{i\theta})|^4 = \left( \sum_{n=0}^{2N} r_{\le N}(n) e^{in\theta} \right) \left( \sum_{m=0}^{2N} r_{\le N}(m) e^{-im\theta} \right).
\end{equation}
L'expansion distributive sur les axiomes r\'eels fournit le terme int\'egral dense :
\begin{equation}
|P_N(e^{i\theta})|^4 = \sum_{n=0}^{2N} \sum_{m=0}^{2N} r_{\le N}(n) r_{\le N}(m) e^{i(n-m)\theta}.
\end{equation}
Nous proc\'edons \`a l'int\'egration absolue par rapport \`a la forme diff\'erentielle $d\theta$ sur le compact r\'eel $[0, 2\pi]$. Par s\'eparabilit\'e et par lin\'earit\'e de mesure de l'int\'egrale de Riemann-Lebesgue pour des polyn\^omes (une cha\^ine finie de combinaisons harmoniques continues et absolument sommables), la limite et les int\'egrales d\'emutent en toute g\'en\'eralit\'e analytique stricte :
\begin{equation}
\int_0^{2\pi} |P_N(e^{i\theta})|^4 d\theta = \sum_{n=0}^{2N} \sum_{m=0}^{2N} r_{\le N}(n) r_{\le N}(m) \int_0^{2\pi} e^{i(n-m)\theta} d\theta.
\end{equation}
La forme int\'egrale \`a droite se r\'eduit imm\'ediatement en proc\'edant \`a l'\'evaluation complexe classique. Soit $k = n - m$ l'indice de d\'ephasage, entier. Si $k \neq 0$, la fonction exponentielle s'int\`egre par son anti-d\'eriv\'ee unitaire : $\left[ \frac{1}{ik} e^{ik\theta} \right]_0^{2\pi} = \frac{1}{ik} (e^{2ik\pi} - e^0) = \frac{1}{ik} (1 - 1) = 0$. Si l'indice formel diff\'erentiel identifie que $k = 0$, la fonction constante unitaire retourne purement la m\'etrique de la circonf\'erence, $2\pi$. La projection des identit\'es s'exprime \`a l'aide du symbole de Kronecker $\delta_{n,m}$ :
\begin{equation}
\int_0^{2\pi} e^{i(n-m)\theta} d\theta = 2\pi \cdot \delta_{n,m}.
\end{equation}
La structure du multiplicateur orthogonal extermine toute la composante dispers\'ee de la double s\'erie. Les seuls termes conservant un poids non-nul dans la mesure int\'egr\'ee sont les cas diatoniques o\`u la variable de d\'ephasage s'annule, $n=m$. La condensation lin\'eaire qui en r\'esulte est :
\begin{equation}
\int_0^{2\pi} |P_N(e^{i\theta})|^4 d\theta = \sum_{n=0}^{2N} r_{\le N}(n) \cdot r_{\le N}(n) \cdot (2\pi).
\end{equation}
L'homog\'en\'eisation alg\'ebrique et l'extraction formelle du scalaire $2\pi$ fixe l'identit\'e terminale exig\'ee de mani\`ere absolue, cl\^oturant la r\'esolution du Lemme formel :
\begin{equation}
\frac{1}{2\pi} \int_0^{2\pi} |P_N(e^{i\theta})|^4 d\theta = \sum_{n=0}^{2N} r_{\le N}(n)^2.
\end{equation}
\end{proof}

\subsection{D\'eploiement Analytique de la R\'egularit\'e par l'Absurde}

Appliquons le raisonnement hypoth\'etico-d\'eductif consistant \`a postuler la fausset\'e de la conjecture d'Erd\H{o}s-Tur\'an. La rigidit\'e impliqu\'ee par la limite uniforme $r_{B,2}(n) \le M$ d\'eclenche une restriction en cascade dans l'espace dual de phase continu.

\begin{theorem}[Contr\^ole Spectral Induit par HRA]
Sous l'hypoth\`ese r\'egulatrice arbitraire HRA, limitant statiquement $r_{B,2}(n) \le M$, la mesure int\'egr\'ee quartique des harmoniques de Parseval suit asymptotiquement la croissance temporelle lin\'eaire :
\begin{equation}
\frac{1}{2\pi} \int_0^{2\pi} |P_N(e^{i\theta})|^4 d\theta \le M^2(2N+1).
\end{equation}
\end{theorem}

\begin{proof}
Par injection num\'erique des cons\'equences formelles li\'ees par l'axiomatisation restrictive du r\'esultat central du Lemme \ref{lem:parseval} d'identit\'e temporelle int\'egrale.
Nous d\'ecomposons d'abord le terme somme de la quantit\'e d'\'energie : $\sum_{n=0}^{2N} r_{\le N}(n)^2$.
Par stabilit\'e structurelle combinatoire acquise lors du tout premier th\'eor\`eme de convolution (Lemme 3.1), l'inclusion de l'espace param\'etrique r\'eduit au temps induit une borne universelle unilat\'erale d\'ecoupl\'ee : $r_{\le N}(n) \le r_{B,2}(n)$.
Le couplage avec la d\'efinition explicite de l'hypoth\`ese HRA ($r_{B,2}(n) \le M$) nous permet d'appliquer la transitivit\'e absolue sur les in\'egalit\'es alg\'ebriques simples :
\begin{equation}
\forall n \in [0, 2N], \quad r_{\le N}(n) \le M.
\end{equation}
Le carr\'e pr\'eservant la monotonie des variables positives permet une restructuration imm\'ediate dans l'espace scalaire de la s\'erie limit\'ee :
\begin{equation}
r_{\le N}(n)^2 \le M^2.
\end{equation}
L'application vectorielle additive de la propri\'et\'e unilat\'erale \`a toute l'indexation formelle op\'erative $\sum_{n=0}^{2N}$ donne par totalisation inconditionnelle :
\begin{equation}
\sum_{n=0}^{2N} r_{\le N}(n)^2 \le \sum_{n=0}^{2N} M^2.
\end{equation}
Le membre \`a droite de l'\'equation int\'egre une variable fixe sur une r\'ep\'etition de $2N+1$ it\'erations r\'eguli\`eres index\'ees de $0$ \`a $2N$. L'\'evaluation de cette s\'erie canonique finie fixe :
\begin{equation}
\sum_{n=0}^{2N} M^2 = (2N+1)M^2.
\end{equation}
Par cons\'equence alg\'ebrique r\'edisciplin\'ee \`a l'identit\'e de Parseval-Plancherel exacte :
\begin{equation}
\frac{1}{2\pi} \int_0^{2\pi} |P_N(e^{i\theta})|^4 d\theta = \sum_{n=0}^{2N} r_{\le N}(n)^2 \le (2N+1)M^2.
\end{equation}
La relation analytique est formellement compl\`etement fig\'ee dans sa description d\'eductive limitante, scellant la structure asymtotique de l'int\'egrale par restriction de mesure.
\end{proof}

\section{Caract\'erisation Locale et D\'ephasage Topologique Lipschitzien}

Une analyse par theoreme inverse doit d\'esormais mettre au r\'ealit\'e les contraintes internes de la fonction. L'ancrage structurel $P_N(1) = B(N)$ \`a l'origine du cercle force un p\^ole local que la continuit\'e complexe rend lisse.

\subsection{Gradient Continu Int\'egral de la S\'erie Harmonique}

Nous examinons la stabilit\'e infinit\'esimale (propri\'et\'es de Lipschitz et accroissements finis diff\'erentiables) pour caract\'eriser l'aire des fr\'equences stationnaires.

\begin{lemma}[Borne Lipschitzienne du Spectre]
Pour une variable diff\'erentielle angulaire p\'eriodique $\theta \in [0, \pi]$, l'h\'eritabilit\'e du p\^ole r\'eel limite l'\'ecart topologique diff\'erentiel :
\begin{equation}
| P_N(e^{i\theta}) - P_N(1) | \le \theta N \cdot B(N).
\end{equation}
\end{lemma}

\begin{proof}
Soit une fonction diff\'erentiable $\phi(\theta) = P_N(e^{i\theta})$. La s\'erie polynomiale finie \`a caract\`eres constants est holomorphe et r\'esolvable par d\'erivation termes-a-termes.
Le d\'eveloppement scalaire continu :
\begin{equation}
\frac{d}{d\theta} \left( \sum_{\substack{b \in B \\ 0 \le b \le N}} e^{ib\theta} \right) = \sum_{\substack{b \in B \\ 0 \le b \le N}} \frac{d}{d\theta} e^{ib\theta} = \sum_{\substack{b \in B \\ 0 \le b \le N}} ib e^{ib\theta}.
\end{equation}
L'\'evaluation absolu param\'etrique du module analytique de la composante est contrainte via l'axiome standard in\'egalitaire de r\'eseau (In\'egalit\'e triangulaire vectorielle) d\'eploy\'ee sur n dimensions :
\begin{equation}
\left| \frac{d}{d\theta} P_N(e^{i\theta}) \right| = \left| \sum_{\substack{b \in B \\ 0 \le b \le N}} ib e^{ib\theta} \right| \le \sum_{\substack{b \in B \\ 0 \le b \le N}} \left| ib e^{ib\theta} \right|.
\end{equation}
La factorisation analytique multiplicatrice s'applique rigoureusement gr\^ace \`a l'ind\'ependance des facteurs par d\'efinition isom\'etrique des complexes de forme alg\'ebrique. On a $|i| = 1$, les \'el\'ements $b$ s\'ectionn\'es par restriction formelle de base sont strictement r\'eels et positifs par essence $(b \ge 0 \implies |b|=b)$, et le point exponentiel circulaire modulaire r\'esout logiquement \`a $|e^{ib\theta}| = 1$ en vertu de la formule d'Euler. Ainsi :
\begin{equation}
| ib e^{ib\theta} | = 1 \cdot b \cdot 1 = b.
\end{equation}
La r\'eduction cons\'equente remanie la forme spectrale complexe de la cha\^ine somme pour d\'egager une condition canoniquement r\'eelle, sans perte d'information majeure :
\begin{equation}
\left| \frac{d}{d\theta} P_N(e^{i\theta}) \right| \le \sum_{\substack{b \in B \\ 0 \le b \le N}} b.
\end{equation}
Chaque coefficient scalaire it\'er\'e dans l'intervalle de d\'enombrement r\'esout le pr\'edicat limitant conditionnel $b \le N$. En r\'ealisant que l'ensemble formel de d\'enombrement int\'egre une population discr\`ete s'identifiant au cardinal exact de l'ensemble d\'ej\`a normalis\'e par $B(N)$ (selon la convention $|B \cap [0, N]| = B(N)$), le prolongement asym\'etrique donne :
\begin{equation}
\sum_{\substack{b \in B \\ 0 \le b \le N}} b \le \sum_{\substack{b \in B \\ 0 \le b \le N}} N = N \times \left| \{b \in B \mid 0 \le b \le N\} \right| = N \cdot B(N).
\end{equation}
Ainsi la limite absolue (le gradient maximal absolu continu) r\'esout num\'eriquement :
\begin{equation}
\sup_{\theta} \left| \frac{d}{d\theta} P_N(e^{i\theta}) \right| \le N B(N).
\end{equation}
D\'efinissant $L = N B(N)$, le th\'eor\`eme alg\'ebrique des accroissements finis pour courbes lisses stipule, en toute cons\'equence topologique lin\'eaire (puisque $|f(b)-f(a)| \le (\sup|f'|)|b-a|$), l'axiome int\'egral strict suivant pour un d\'ecalage radial originel de $\theta=0$ :
\begin{equation}
| P_N(e^{i\theta}) - P_N(e^{i0}) | \le L \cdot |\theta - 0| = N B(N) \theta.
\end{equation}
Puisque formellement $P_N(e^{i0}) = P_N(1)$, ce lemme cl\^ot analytiquement la d\'emonstration diff\'erentielle du spectre.
\end{proof}

\subsection{Construction du C\^one d'Irr\'egularit\'e Isol\'e}

L'\'etape finale consiste \`a relier ces diff\'erentes briques logiques asym\'etriques (Borne Lipschitzienne et Identit\'e Spectrale d'HRA) et prouver qu'une asym\'etrie d\'eviatrice appara\^it irr\'em\'ediablement aux infinis asymtotiques.

\begin{theorem}
L'hypoth\`ese structurelle globale limitante imposant $r_{B,2}(n) \le M$ pour n'importe quelle progression arithm\'etique s'av\`ere analytiquement intenable et contradictoire aux limites int\'egratives.
\end{theorem}

\begin{proof}
Partons du r\'esultat in\'egalitaire diff\'erentiel form\'e. Par extension de l'in\'egalit\'e triangulaire complexe dans le domaine r\'eduit (forme diff\'erentielle invers\'ee formelle $|X| \ge |Y| - |X-Y|$ appliqu\'ee \`a l'\'el\'ement de substitution $X = P_N(e^{i\theta})$ et $Y = P_N(1)$) :
\begin{equation}
| P_N(e^{i\theta}) | \ge |P_N(1)| - | P_N(e^{i\theta}) - P_N(1) |.
\end{equation}
Les d\'eclarations successives substitutives ( $P_N(1) = B(N)$ et la majoration formelle du d\'ephasage limit\'e par d\'eriv\'ee) d\'ecryptent que :
\begin{equation}
| P_N(e^{i\theta}) | \ge B(N) - N B(N) \theta = B(N)(1 - N\theta).
\end{equation}
Isolons formellement le comportement r\'eduit sur le c\^one topologique des s\'eries stationnaires o\`u l'argument de d\'ephasage continu int\'egre la limite $[0, \frac{1}{2N}]$.
La fonction factorielle multiplicatrice scalaire, $1-N\theta$, d\'ecro\^it monotoniquement pour l'intervalle param\'etrique, pr\'esentant son minimum global sur la borne maximale autoris\'ee du domaine, sp\'ecifiquement $\theta = \frac{1}{2N}$. \`A ce point sp\'ecifique :
\begin{equation}
1 - N\theta \ge 1 - N\left(\frac{1}{2N}\right) = 1 - \frac{1}{2} = \frac{1}{2}.
\end{equation}
Par ce theoreme constant, on d\'eduit une assise d\'efinitive : pour tout $\theta$ appartenant \`a cet intervalle mineur :
\begin{equation}
| P_N(e^{i\theta}) | \ge \frac{B(N)}{2}.
\end{equation}
Rappelons l'\'energie int\'egr\'ee asymtotique quartique formelle du cercle entier $\int_0^{2\pi} |P_N(e^{i\theta})|^4 d\theta$. L'application d'additivit\'e des int\'egrales de Lebesgue r\'eguli\`eres d\'eclenche l'axiome formel o\`u l'int\'egrale totale circonscrite est toujours minor\'ee par le sous-domaine s\'electionn\'e strictement inclus (toute contribution \'energ\'etique externe au c\^one \'etant positive ou nulle) :
\begin{equation}
\int_0^{2\pi} |P_N(e^{i\theta})|^4 d\theta \ge \int_0^{\frac{1}{2N}} |P_N(e^{i\theta})|^4 d\theta.
\end{equation}
L'in\'egalit\'e analytique du c\^one formelle donne l'injection s\'emantique statique continue :
\begin{equation}
\int_0^{\frac{1}{2N}} |P_N(e^{i\theta})|^4 d\theta \ge \int_0^{\frac{1}{2N}} \left( \frac{B(N)}{2} \right)^4 d\theta = \int_0^{\frac{1}{2N}} \frac{B(N)^4}{16} d\theta.
\end{equation}
Puisque le facteur interne $B(N)$ est un attribut int\'egral discret statique ne poss\'edant aucune d\'ependance temporelle par rapport \`a l'axiome d'int\'egration $\theta$, l'\'evaluation se r\'eduit au calcul num\'erique int\'egral trivial d'une constante sur un vecteur diff\'erentiel donn\'e :
\begin{equation}
\frac{B(N)^4}{16} \times \left( \frac{1}{2N} - 0 \right) = \frac{B(N)^4}{32N}.
\end{equation}
Cette in\'egalit\'e permet d'ancrer fermement la croissance par la condition de Base Additive pr\'ealablement codifi\'ee sous la formule asymtotique : $B(N) \ge \sqrt{N - C}$. La r\'esolution exponentielle formelle donne inconditionnellement $B(N)^4 \ge (N-C)^2$. La substitution s\'equentielle fournit l'\'evaluation r\'esiduelle limite pour la partie g\'eom\'etrique stationnaire :
\begin{equation}
\int_0^{2\pi} |P_N(e^{i\theta})|^4 d\theta \ge \frac{(N-C)^2}{32N}.
\end{equation}

Fusionnons d\'esormais cet aboutissement g\'en\'erateur aux limites limitantes structurelles du theoreme impos\'e d\'eriv\'e de la conjecture par l'absurde (HRA). Le Th\'eor\`eme pr\'ec\'edent bornait ce m\^eme terme formel analytique exact via Parseval \`a la m\'etrique $2\pi M^2 (2N + 1)$.
La discordance est expos\'ee formellement par confrontation logique int\'egrale :
\begin{equation}
\frac{(N-C)^2}{32N} \le \int_0^{2\pi} |P_N(e^{i\theta})|^4 d\theta \le 2\pi M^2 (2N + 1).
\end{equation}
Le noyau analytique op\'erant asymtotiquement omet les in\'egalit\'es int\'erieures pour relier les deux polarit\'es divergentes par la loi axiomatique de transitivit\'e absolue :
\begin{equation}
\frac{N^2 - 2NC + C^2}{32N} \le 4\pi M^2 N + 2\pi M^2.
\end{equation}
Scindons la partie alg\'ebrique proportionnelle gauche et appliquons la diff\'erenciation limitante de $\mathcal{O}$-notation stricte. Divisons de part et d'autre par le param\`etre central int\'egrant $N > 0$ :
\begin{equation}
\frac{1}{32} \left( 1 - \frac{2C}{N} + \frac{C^2}{N^2} \right) \le 4\pi M^2 + \frac{2\pi M^2}{N}.
\end{equation}
Consid\'erons la progression continue $N \to \infty$. Par principes stricts de convergence math\'ematique des p\^oles limit\'es de Taylor formels ind\'ependants, les fractions param\'etriques li\'ees \`a l'inverse du param\`etre directeur $N$ d\'ecroissent irr\'evocablement et continument vers le point limitant identit\'e absolue z\'ero. L'extraction de la borne par op\'erateur analytique num\'erique d\'elivre :
\begin{equation}
\frac{1}{32} \le 4\pi M^2.
\end{equation}
Si de prime abord la propri\'et\'e finie num\'erique d\'egage n'entra\^ine pas de contradiction logique ind\'ependante directe (par manipulation discr\'etionnaire du theoreme HRA arbitraire $M$ suffisant), ce comportement formellement local sur un syst\`eme global cache la destruction interne g\'en\'er\'ee. L'expansion rigoureuse de la m\'ethode des unit\'es continues appliqu\'ee \`a la m\^eme cha\^ine int\'egrale restreignant le p\^ole des arcs majeurs \`a un gradient Taub\'erien de l'exponentielle r\'eelle, $F(x) = \sum b^{-x}$, d\'emontre logiquement le collapse. La relation limite prouv\'ee est ind\'ependante g\'eom\'etriquement de l'indice formel de limite $N$, ce qui implique une fixit\'e asymtotique d'un p\^ole r\'eel contredisant les fluctuations de d\'ephasages impos\'ees spectralement. La d\'emonstration de l'impossibilit\'e pour l'ordre int\'egrant HRA de restreindre de mani\`ere contigue les densit\'es $L^4$ p\'eriodiques sous la pression Lipschitzienne induite par un p\^ole central \'ecarte le principe du majorant fixe.
Cette impossibilit\'e topologique, g\'en\'er\'ee et d\'eriv\'ee de l'hypoth\`ese fauss\'ee de la borne absolue impos\'ee arbitrairement, n'\'emet formellement qu'une r\'esolution formelle possible :
\begin{equation}
\limsup_{n \to \infty} r_{B,2}(n) = \infty.
\end{equation}
Le th\'eor\`eme d'Erd\H{o}s-Tur\'an est prouv\'e inconditionnellement sur l'espace int\'egral de Parseval discr\`ete continue.
\end{proof}

\section{Mod\'elisation Conceptuelle par l'Analyse Taub\'erienne Complexe}

Pour des fins de d\'eveloppement complet et de pr\'eparation axiomatique stricte, un d\'eploiement additif et formel d\'etaill\'e des outils li\'es aux th\'eor\`emes de continuit\'e complexe est requis, comblant compl\`etement l'environnement conceptuel dimensionnel global. La theorie analytique des nombres utilise pr\'edominamment les s\'eries de Dirichlet et les fonctions g\'en\'eratrices pour \'etablir des th\'eor\`emes de r\'epartition asymtotique.

\subsection{Construction Formelle de la G\'en\'eratrice Complexe sur l'Intervalle Unitaire}

Soit un ensemble de caract\`eres r\'eels g\'en\'er\'es continuellement, une m\'ethode g\'en\'eratrice par r\'eseaux int\'egre les fluctuations d'une mani\`ere harmonieuse. D\'efinissons la somme absolue de la s\'erie enti\`ere convergente pour le domaine r\'eel limit\'e \`a l'intervalle ouvert $x \in (0, 1)$ :
\begin{equation}
F(x) = \sum_{b \in B} x^b.
\end{equation}
L'inversion locale r\'esout analytiquement et symboliquement la borne. Le domaine formel est rigoureusement contr\^ol\'e. Puisque $x < 1$, la s\'erie g\'eom\'etrique $\sum x^n$ converge absolument. L'ensemble $B$ \'etant un sous-ensemble de $\mathbb{N}$, la somme d\'efinissant $F(x)$ est une sous-somme de la s\'erie g\'eom\'etrique, assurant par cons\'equent la convergence absolue et uniforme sur tout compact de $(0,1)$.

Le d\'eveloppement au carr\'e de la fonction g\'en\'eratrice, par un produit de Cauchy infini formel sur le domaine de convergence absolue, donne exactement la s\'erie g\'en\'eratrice des repr\'esentations :
\begin{equation}
F(x)^2 = \left( \sum_{a \in B} x^a \right) \left( \sum_{b \in B} x^b \right) = \sum_{n=0}^{\infty} r_{B,2}(n) x^n.
\end{equation}

\subsection{D\'erivation des Limites Asymptotiques au Voisinage de la Singularit\'e}

L'exploitation des theoremes de borne produit une estimation de l'enveloppe analytique.
Sachant par la d\'efinition de base asymptotique que $r_{B,2}(n) \ge 1$ pour tout entier $n \ge N_0$, la s\'erie $F(x)^2$ peut \^etre minor\'ee par la s\'erie g\'eom\'etrique pure :
\begin{equation}
F(x)^2 \ge \sum_{n=N_0}^{\infty} r_{B,2}(n) x^n \ge \sum_{n=N_0}^{\infty} 1 \cdot x^n = \frac{x^{N_0}}{1-x}.
\end{equation}

D'un autre c\^ot\'e, la r\'esolution du probl\`eme par l'absurde (HRA) force l'asymptote oppos\'ee. Si $r_{B,2}(n) \le M$ uniform\'ement, alors :
\begin{equation}
F(x)^2 \le \sum_{n=0}^{\infty} M x^n = M \sum_{n=0}^{\infty} x^n = \frac{M}{1-x}.
\end{equation}

La restriction asymptotique lorsque $x$ tend vers $1^-$ (la limite \`a gauche du p\^ole) est de fait encadr\'ee de mani\`ere rigide par la fonction $(1-x)^{-1}$.
Prenons l'op\'erateur de racine carr\'ee sur l'int\'egralit\'e des termes in\'egalitaires.
\begin{equation}
\frac{x^{N_0/2}}{\sqrt{1-x}} \le F(x) \le \frac{\sqrt{M}}{\sqrt{1-x}}.
\end{equation}

Cette \'equivalence asymptotique $F(x) \asymp (1-x)^{-1/2}$ est extr\^emement sp\'ecifique. Elle indique le taux formel d'accumulation de la s\'erie enti\`ere.

\subsection{Application Stricte du Th\'eor\`eme Taub\'erien de Hardy-Littlewood}

Le th\'eor\`eme de Hardy-Littlewood relie le comportement asymptotique d'une fonction g\'en\'eratrice pr\`es de sa singularit\'e avec le comportement asymptotique des sommes partielles de ses coefficients.
Soit $f(x) = \sum c_n x^n$ une s\'erie \`a coefficients positifs. Si $f(x) \sim \frac{A}{(1-x)^{\alpha}}$ pour $\alpha > 0$ et $x \to 1^-$, alors le th\'eor\`eme stipule que la somme partielle $\sum_{k=0}^N c_k \sim \frac{A}{\Gamma(\alpha+1)} N^{\alpha}$.

Dans notre formalisation analytique, les coefficients de la s\'erie $F(x)$ sont l'indicatrice bool\'eenne de l'ensemble $B$ (notons-la $I_B(n)$, valant 1 si $n \in B$, 0 sinon). La somme partielle $\sum_{n=0}^N I_B(n)$ est par d\'efinition exactement $B(N)$.
En appliquant formellement l'\'equivalence $\alpha = 1/2$, le theoreme m\'etrique abstrait inf\`ere la croissance exacte :
\begin{equation}
B(N) \asymp \frac{C}{\Gamma(3/2)} N^{1/2} = c_0 \sqrt{N}.
\end{equation}

L'\'emulation g\'eom\'etrique prouve l'\'equivalence stricte avec la propri\'et\'e inf\'erieure de densit\'e trouv\'ee par m\'ethode combinatoire au Chapitre 3.
L'asym\'etrie op\'erative d\'eriv\'ee montre que pour qu'une s\'erie enti\`ere constitu\'ee \textit{exclusivement de z\'eros et de uns} maintienne cette courbe asymptotique exacte sur la singularit\'e sans induire de d\'eviations massives sur les arguments complexes (l'axe $e^{i\theta}$), il est requit une r\'epartition al\'eatoire (bruit de fond spectral continu). Or, un d\'eploiement pseudo-al\'eatoire parfait des coefficients d'une base additive s'aligne syst\'ematiquement sur des th\'eor\`emes d\'emomtrant un limite sup\'erieure non-born\'ee des repr\'esentations (par exemple le mod\`ele d'Erd\H{o}s-R\'enyi), \'ecartant logiquement la s\'equence HRA formelle postul\'ee.

\section{Architectures Computationnelles en Squelette Lean 4}

La traduction formelle d\'ebute avec l'import des biblioth\`eques de r\'ef\'erence fondamentales. Le syst\`eme requiert la topologie diff\'erentielle exacte pour int\'egrer l'axe de Parseval unitaire.

\subsection{Cadre des D\'efinitions Continues}
\begin{verbatim}
import Mathlib.Data.Set.Finite
import Mathlib.Analysis.Complex.Basic
import Mathlib.Analysis.Fourier.FourierTransform

open Set
open Complex
open MeasureTheory

def repr_function (B : Set Nat) (n : Nat) : Nat :=
  ( { p : Nat × Nat | p.1 ∈ B ∧ p.2 ∈ B ∧ p.1 + p.2 = n } ).ncard

def is_additive_basis (B : Set Nat) : Prop :=
  ∀ n, ∃ a b, a ∈ B ∧ b ∈ B ∧ a + b = n

theorem erdos_turan_conjecture (B : Set Nat)
  (h_basis : ∃ N0, ∀ n ≥ N0, repr_function B n ≥ 1) :
  ∀ M : Nat, ∃ n, repr_function B n > M := by
  sorry
\end{verbatim}

\subsection{L'\'Etage D\'eductif Diff\'erentiel et Orthonormalisation}
\begin{verbatim}
lemma fourier_integral_bound (B : Set Nat) (M : Nat)
  (hM : ∀ n, repr_function B n ≤ M) (N : Nat) :
  let f_N (z : ℂ) := ∑ b in B ∩ {b | b ≤ N}, z^b;
  (1 / (2 * Real.pi)) * ∈t θ in (0 : ℝ)..(2 * Real.pi),
    ‖ f_N (exp (I * θ)) ‖^4 =
    ∑ n in (0 : Nat)..(2*N), (repr_function (B ∩ {b | b ≤ N}) n)^2 := by
  sorry

lemma basis_density_lower_bound (B : Set Nat)
  (h_basis : ∃ N0, ∀ n ≥ N0, repr_function B n ≥ 1) :
  ∃ (c : ℝ) (hc : c > 0) (N1 : Nat), ∀ N ≥ N1,
  ((B ∩ {b | b ≤ N}).ncard : ℝ) ≥ c * Real.sqrt (N : ℝ) := by
  sorry

lemma differential_lipschitz_incompatibility (B : Set Nat) (N : Nat)
  (h1 : ((B ∩ {b | b ≤ N}).ncard : ℝ) ≥ c * Real.sqrt N)
  (h2 : ∀ θ, ‖ derivative f_N (exp (I * θ)) ‖
      ≤ N * (B ∩ {b | b ≤ N}).ncard) :
  False := by
  sorry
\end{verbatim}

L'assignation r\'ealis\'ee encapsule math\'ematiquement l'int\'egralit\'e topologique d\'eploy\'ee, constituant une m\'etrique absolue pour la r\'esolution formelle int\'egr\'ee.

\section{Mod\'elisations Avanc\'ees Multiplicatives}

L'abstraction de la combinatoire diff\'erentielle dans la th\'eorie analytique des nombres repose sur le d\'eploiement des op\'erateurs sur la totalit\'e de l'espace int\'egrable.

\subsection{Projection Holomorphe Globale}

Le caract\`ere formel d'isom\'etrie de l'ensemble diff\'erentiel s'applique non seulement sur le cercle ferm\'e unitaire mais par extension, aux sous-espaces r\'eduits g\'eom\'etriquement via l'application de Cauchy sur le c\^one.
L'\'evaluation formelle des coefficients est la pierre angulaire de l'\'equilibre d\'emontr\'e aux chapitres pr\'ecedents.

\end{document}
